\begin{apartats}

\apartat[1,25]{0,75}
Calcula les integrals $\displaystyle\int\left(e^{2x}+2^x\right)dx$ i
$\displaystyle\int\frac{4}{2x+3}\,dx$.

\begin{solucio}
$\displaystyle\int\left(e^{2x}+2^x\right)dx=\frac{e^{2x}}{2}+\frac{2^x}{\ln2}+k$. Com que la derivada de
$2x+3$ és 2, $\displaystyle\int\frac{4}{2x+3}\,dx=2\int\frac{2}{2x+3}\,dx=2\ln|2x+3|+k$.
\end{solucio}

\begin{tria}{immediates-altres}
\itemtria{trigonometriques}{1}{1,25}
Calcula $\displaystyle\int\bigl(\cos(3x)+3\sin x\bigr)\,dx$ i $\displaystyle\int\sin(x-\pi)\,dx$.
Comprova els resultats derivant.

\begin{solucio}
$\displaystyle\int\bigl(\cos(3x)+3\sin x\bigr)\,dx=\frac{\sin(3x)}{3}-3\cos x+k$: la derivada és
$\frac{3\cos(3x)}{3}+3\sin x$.\\
$\displaystyle\int\sin(x-\pi)\,dx=-\cos(x-\pi)+k$: la derivada és $\sin(x-\pi)$.
\end{solucio}

\itemtria{quocients}{1}{1,25}
Calcula $\displaystyle\int\frac{x^2+1}{x}\,dx$ i $\displaystyle\int\frac{(x+1)^2}{\sqrt x}\,dx$,
separant cada quocient en una suma.

\begin{solucio}
$\frac{x^2+1}{x}=x+\frac1x$, i $\displaystyle\int\frac{x^2+1}{x}\,dx=\frac{x^2}{2}+\ln|x|+k$.\\
$\frac{(x+1)^2}{\sqrt x}=\frac{x^2+2x+1}{x^{1/2}}=x^{3/2}+2x^{1/2}+x^{-1/2}$, i
\[
\int\frac{(x+1)^2}{\sqrt x}\,dx=\frac25x^{5/2}+\frac43x^{3/2}+2x^{1/2}+k
=\frac25x^2\sqrt x+\frac43x\sqrt x+2\sqrt x+k.
\]
\end{solucio}
\end{tria}

\begin{nomesllarg}
\apartat{0,75}
Calcula $\displaystyle\int\left(\frac{1}{1+x^2}+\frac{1}{\sqrt{1-x^2}}\right)dx$ i comprova el
resultat derivant.

\begin{solucio}
Són dues integrals immediates: $\operatorname{arctg}x+\arcsin x+k$. Derivant,
$\frac{1}{1+x^2}+\frac{1}{\sqrt{1-x^2}}$, que és la funció de partida.
\end{solucio}
\end{nomesllarg}

\end{apartats}
